\chapter{Objectives and Status}\label{chap:8}

\begin{quote}
\emph{``Tell me how you measure me, and I'll tell you how I'll
behave.''} --- Eliyahu Goldratt
\end{quote}

This chapter covers the objective on the model side, and the status
code, bounds, and parameters that let us interpret the result on the
solver side.

\section{Stating the Objective}\label{sec:8.1}

A CP-SAT model has at most one objective.

\begin{lstlisting}[language=Python]
model.minimize(linear_expression)
model.maximize(linear_expression)
\end{lstlisting}

The expression must be \emph{linear} in the variables: sums, scaled
sums, and integer-coefficient combinations. As elsewhere in CP-SAT,
mixing integer and Boolean variables is fine: a Boolean is a
\(\{0, 1\}\) integer in arithmetic, and negations via
\lstinline!~! are accepted in objectives just like in
constraints.

\begin{lstlisting}[language=Python]
# Total cost of selected items
model.minimize(sum(cost[i] * take[i] for i in range(n)))

# Mixing terms — Booleans contribute 0 or 1, with negation if useful
model.minimize(sum(weight[i] * x[i] for i in range(n)) - 100 * ~all_done)
\end{lstlisting}

Calling \lstinline!minimize! or
\lstinline!maximize! a second time \emph{replaces} the
previous objective; it does not add to it. CP-SAT does not have a
built-in concept of multiple simultaneous objectives. To combine several
goals (\emph{minimize lateness AND minimize cost}), the modeler
commits to a combination strategy, and \Cref{sec:8.3} covers the two standard
ones.

A model with no objective is a \emph{constraint-satisfaction} problem:
the solver returns any feasible solution it can find, with no preference
among them. CP-SAT is strong in this mode: Sudoku, schedule
feasibility, and configuration validation all fall here. A dummy
objective like \lstinline!model.minimize(0)! is not
needed; feasibility-only tells the solver \emph{find any solution},
which is the natural question for many problems and frees the search
from optimization overhead.

\section{Nonlinear Objectives via Auxiliaries}\label{sec:8.2}

The objective expression must be linear, but \emph{the quantity we want
to optimize} often is not: a maximum, an absolute value, a
multiplicative cost. The standard move: introduce an auxiliary variable
bound to the nonlinear quantity, then make the objective linear in the
auxiliary. The pattern is the same one \Cref{chap:7}'s equality helpers
(\Cref{sec:7.10}) used for nonlinear \emph{constraints}, applied to the
\emph{objective}.

The most common case is \emph{minimize a maximum}: a makespan, a
worst-case latency, a peak load. The direct route uses
\lstinline!add_max_equality!:

\begin{lstlisting}[language=Python]
makespan = model.new_int_var(0, horizon, "makespan")
model.add_max_equality(makespan, [end[t] for t in tasks])
model.minimize(makespan)
\end{lstlisting}

The linear-bound trick from \Cref{sec:7.10} expresses the same thing without the
equality constraint:

\begin{lstlisting}[language=Python]
makespan = model.new_int_var(0, horizon, "makespan")
for t in tasks:
    model.add(makespan >= end[t])
model.minimize(makespan)
\end{lstlisting}

The bounds say \emph{makespan is at least every finish time}, and the
minimization pulls the bound tight to the actual maximum on its own. The
mirror pattern (\emph{maximize a minimum}) uses
\(m \le \text{value}\) for each value with
\lstinline!model.maximize(m)!.

For other nonlinear shapes, the auxiliary-variable pattern carries
through unchanged:

\begin{lstlisting}[language=Python]
# Total absolute deviation from per-task target dates
total_dev = model.new_int_var(0, big, "total_dev")
deviations = []
for t in tasks:
    d = model.new_int_var(0, big, f"dev_{t}")
    model.add_abs_equality(d, finish_time[t] - target[t])
    deviations.append(d)
model.add(total_dev == sum(deviations))
model.minimize(total_dev)
\end{lstlisting}

Same recipe: scaffold an auxiliary for each nonlinear piece (here an
absolute value per task), sum the auxiliaries, make the sum the
objective. The pattern handles modulo objectives, multiplicative
penalties, and any other nonlinear corner of integer arithmetic.

\section{Combining Multiple Goals}\label{sec:8.3}

When the model has more than one goal, two standard patterns combine
them into a single objective. They differ in whether the goals are
negotiable against each other (weighted sum) or strictly prioritized
(lexicographic).

\subsection{Weighted sum}\label{sec:8.3.1}

The simplest combination: pick weights and minimize the weighted sum.

\begin{lstlisting}[language=Python]
W_COST       = 1
W_LATENESS   = 100
W_PREFERENCE = 10

model.minimize(
      W_COST       * total_cost
    + W_LATENESS   * total_lateness
    - W_PREFERENCE * preference_score
)
\end{lstlisting}

The weights encode how much one unit of one goal is worth in units of
another: \emph{one minute of lateness is worth one hundred dollars of
cost}. Weighted sums are the right tool when the goals are genuinely
commensurable and the trade-off rate is known or can be estimated.

Setting one weight extremely large to mean ``this matters most''
(\lstinline!W_LATENESS = 10**9! or similar) tends to
backfire: the objective becomes effectively \emph{only} the
high-weighted term with the others lost in the noise, and the solver can
struggle numerically when coefficients span many orders of magnitude.
When a goal is genuinely top-priority (optimized first, others as
tie-breakers), \emph{lexicographic ordering} expresses that directly.

\subsection{Lexicographic priorities}\label{sec:8.3.2}

When goals are strictly ranked (\emph{first satisfy lateness, then
minimize cost}), solve in \emph{stages}. Each stage optimizes one
objective; before moving to the next, the previous objective's optimal
value is locked in as a constraint.

\begin{lstlisting}[language=Python]
# Stage 1: minimize lateness
model.minimize(total_lateness)
status = solver.solve(model)
best_lateness = round(solver.objective_value)

# Stage 2: lock lateness in, optimize cost as a tie-breaker
model.add(total_lateness <= best_lateness)
# model.clear_objective() is the explicit form; minimize() replaces it.
model.minimize(total_cost)

# Warm-start: the stage-1 solution is still feasible and usually
# near-optimal for the new objective — hint from it.
for v in decision_vars:
    model.add_hint(v, solver.value(v))
solver.parameters.repair_hint = True

status = solver.solve(model)
\end{lstlisting}

The pattern generalizes to any number of stages: optimize, lock,
optimize the next. The result is a strict priority ordering: every
later objective is optimized only over the optimal-for-the-earlier set.
The hinting step matters: changing the objective does not change
feasibility, so the previous solution is a free starting point for the
next stage and can speed up the next solve substantially.

A practical note: by the time a model reaches its third or fourth
lexicographic stage, the search space has been trimmed aggressively by
the locked-in constraints, and the remaining problem can be unexpectedly
hard. Setting
\lstinline!solver.parameters.max_time_in_seconds! on
the secondary stages keeps a tie-breaker from consuming the entire time
budget.

\section{Reading the Solver's Answer}\label{sec:8.4}

The solver returns a \emph{status}, an \emph{objective value} (when an
objective is set), a \emph{best objective bound}, and a wall-clock time.

\subsection{Solver parameters}\label{sec:8.4.1}

A handful of parameters get set on essentially every nontrivial solve:

\begin{lstlisting}[language=Python]
solver = cp_model.CpSolver()
solver.parameters.max_time_in_seconds = 30.0  # soft time limit
solver.parameters.num_workers = 8             # parallel portfolio
solver.parameters.log_search_progress = True  # readable progress log

status = solver.solve(model)
\end{lstlisting}

\lstinline!max_time_in_seconds! is the most important.
Without it, the solver runs until it finds the optimum, proves
infeasibility, or exhausts memory, which on hard instances can mean
\emph{days}. Set a budget appropriate to the problem and the moment.

\lstinline!num_workers! controls the parallel portfolio:
different workers run different strategies and share learned clauses
with each other. Eight is a sensible default on a modern laptop, larger
on a workstation; \Cref{sec:5.3}'s hardware sketch translates directly to a worker
count. The default is \emph{all} cores including hyperthreads, which is
not always best: on heavily contended machines, capping at the number
of physical cores often runs faster.

\lstinline!log_search_progress! turns on a
human-readable trace of the search. When debugging a slow model, the log
is the first place to look: it shows when bounds tighten, when the LP
relaxation bites, and when LNS finds an improvement. The full parameter
list is large; the OR-Tools documentation is the authoritative reference
for the corners.

\subsection{Status codes}\label{sec:8.4.2}

The status returned by \lstinline!solver.solve(...)! is
one of five values, each with a precise meaning:

\begin{symboltable}{lL}
Code & Meaning \\ \midrule
\lstinline!OPTIMAL! & A solution was found \emph{and} proven to be the best possible. \\
\lstinline!FEASIBLE! & A solution was found, but optimality has not been proven. \\
\lstinline!INFEASIBLE! & The solver proved that \emph{no} solution exists. \\
\lstinline!MODEL_INVALID! & The model itself is malformed (empty domains, integer overflow, API misuse). \\
\lstinline!UNKNOWN! & The solver gave up before deciding either way. \\
\end{symboltable}


Three of these are \emph{proofs}: \lstinline!OPTIMAL!,
\lstinline!INFEASIBLE!, and
\lstinline!MODEL_INVALID! are mathematical certainties
about the model. \lstinline!FEASIBLE! is a \emph{partial}
result: there is a solution, and we have it, but we have not proven
that no better one exists. \lstinline!UNKNOWN! is the
absence of any conclusion at all.

The most common confusion is between \lstinline!OPTIMAL!
and \lstinline!FEASIBLE!. A solver that returns
\lstinline!FEASIBLE! after a 30-second time limit has not
failed: it has found a solution and reports its objective value, but
the search ran out of time before proving optimality. The objective
value is the best-known; the next subsection describes how to assess how
close that best-known might be to the actual optimum.

A typical branching pattern:

\begin{lstlisting}[language=Python]
if status in (cp_model.OPTIMAL, cp_model.FEASIBLE):
    # We have a solution we can act on. Read it.
    ...
elif status == cp_model.INFEASIBLE:
    # The problem is over-constrained — the model and the data
    # together describe an impossible instance.
    ...
elif status == cp_model.MODEL_INVALID:
    # The model itself is broken. Read the log carefully.
    ...
else:  # UNKNOWN
    # The solver gave up. More time, more workers, or a tighter model.
    ...
\end{lstlisting}

\lstinline!OPTIMAL! and \lstinline!FEASIBLE!
go to the same branch because the \emph{answer} is usable in both cases:
the difference matters for monitoring and reporting, not for acting
on the result.

\subsection{The bound: what optimality looks like}\label{sec:8.4.3}

When an objective is set, the solver maintains two values. The
\emph{best known objective value} is the best feasible solution found so
far. The \emph{best objective bound} is a proven limit beyond which the
objective cannot improve: for a minimization problem, a \emph{lower
bound} on the optimum; for maximization, an \emph{upper bound}.

\begin{lstlisting}[language=Python]
print(solver.objective_value)       # best-known solution (float)
print(solver.best_objective_bound)  # best proven bound (dual side)
print(solver.wall_time)             # seconds elapsed
\end{lstlisting}

When status is \lstinline!OPTIMAL!,
\lstinline!objective_value! and
\lstinline!best_objective_bound! are equal (modulo
floating-point representation): the solver has both found a solution and
proven that nothing better exists. When status is
\lstinline!FEASIBLE!, the two differ, and the \emph{gap}
between them is what tells us how close the best-known is to provable
optimality.

The standard relative-gap metric:

\[
\text{gap} \;=\; \frac{|\text{objective\_value} - \text{best\_objective\_bound}|}{|\text{objective\_value}|}.
\]

A gap of zero is a proof of optimality. A small gap (a few percent)
means the best-known is within that fraction of the optimum, and
continuing the search may not improve the answer much. A large gap means
the bound is loose, the answer might be far from optimal, and either
more time, a better model, or a stronger formulation is warranted. (The
formula divides by \lstinline!|objective_value|! and is
therefore undefined when the objective evaluates to zero; the OR-Tools
logs handle this by guarding the denominator, but in user code one
usually special-cases the zero-objective case.)

If ``close enough'' is good enough for the application, the solver can
stop on its own once the gap is below a threshold:

\begin{lstlisting}[language=Python]
solver.parameters.relative_gap_limit = 0.01    # stop at 1 %
solver.parameters.absolute_gap_limit = 5       # or within 5 units
\end{lstlisting}

This is often a better fit for production than a fixed time limit,
because it gives the solver permission to stop early on easy instances
and keep working on hard ones.

\lstinline!objective_value! is returned as a Python
\lstinline!float! even on integer-only models. Round it
explicitly if the downstream code expects an integer.

\subsection{What to do with each status}\label{sec:8.4.4}

A short field guide that combines the status, the bound, and the action:

\begin{itemize}
\item
  \textbf{\lstinline!OPTIMAL!}: done. Read the solution
  and act on it.
\item
  \textbf{\lstinline!FEASIBLE!} with a small gap:
  usually done. Read the solution; consider whether the gap matters for
  the application.
\item
  \textbf{\lstinline!FEASIBLE!} with a large gap: the
  model needs more time, a better formulation, or revisiting
  representation choices (\Cref{sec:5.5}) or the model's scope (\Cref{chap:9}).
\item
  \textbf{\lstinline!INFEASIBLE!}: the problem and the
  data, \emph{together}, describe an impossible instance. Use the
  verifier from \Cref{chap:4} to confirm: hand-craft a solution you believe
  should be valid and run it through the evaluator. If the evaluator
  accepts it, the encoding is wrong; if the evaluator rejects it, your
  understanding of the rules was wrong. Common fixes: relax a hard
  constraint into the slack-and-penalty form of \Cref{sec:8.5}, fix a data error,
  or accept that no plan exists for the given inputs.
\item
  \textbf{\lstinline!MODEL_INVALID!}: the model itself
  is broken. Read the log; look for variables with empty domains,
  expressions outside the supported integer range, or API misuses.
  (Logical contradictions between \emph{constraints} surface as
  \lstinline!INFEASIBLE!, not
  \lstinline!MODEL_INVALID!: the latter is reserved
  for structural problems with the model itself.)
\item
  \textbf{\lstinline!UNKNOWN!}: the solver had no time
  to conclude. More time, more workers, or a smaller model.
\end{itemize}

\section{Soft Constraints: Slack and Penalty}\label{sec:8.5}

Some rules are preferred but not absolute: \emph{total hours should
not exceed the cap, but overtime is allowed at a cost}; \emph{the
schedule should respect preferences, but the schedule must be filled
even if some preferences cannot be honored}. The standard CP-SAT pattern
combines a relaxed constraint with a penalty in the objective:

\begin{lstlisting}[language=Python]
overtime = model.new_int_var(0, max_overtime, "overtime")

# The hard cap is relaxed; overtime is the slack.
model.add(total_hours <= regular_cap + overtime)

# The penalty discourages — but does not forbid — overtime.
model.minimize(total_cost + 1000 * overtime)
\end{lstlisting}

The slack variable \lstinline!overtime! measures \emph{how
much} the original cap would have been violated, and its coefficient in
the objective is the per-unit penalty. The solver chooses on its own
when paying for slack is cheaper than the alternatives: the
optimization performs the trade-off.

This is the CP-SAT realization of the \emph{hard versus soft}
distinction from \Cref{sec:4.3}. A \emph{hard} rule is encoded as a constraint
with no slack; a \emph{soft} rule is encoded as a constraint \emph{with}
slack, and the slack's coefficient in the objective is the rule's
per-unit penalty. A model with only hard rules has a binary feasibility
verdict; a model with soft rules has a richer answer that captures
\emph{how much} and \emph{where} the rules were violated.

On the verifier side, the matching \lstinline!Rule!
reports its violation through \lstinline!Metric.cost!
rather than \lstinline!Metric.violation!: slack
semantics stay in lockstep on both sides.

The slack pattern also has a workflow purpose: it is the standard
response to an \lstinline!INFEASIBLE! status from \Cref{sec:8.4}
when the infeasibility is real and the rule is genuinely soft. Rather
than refusing to plan, the model returns the \emph{cheapest} violation
of the soft rule that the rest of the constraints permit. This is often
what is wanted in production: a binary ``no plan exists'' answer is
rarely actionable, while ``the cheapest plan violates rule X by 3 units
at cost 3000'' is.
